% ============================================================
% KAPITEL 1 — EINLEITUNG
% ============================================================

\chapter{Einleitung}
\label{ch:einleitung}

% ------------------------------------------------------------
\section{Motivation und Problemstellung}
\label{sec:ein-motivation}

Ein Nutzer steht in einer virtuellen Szene und gibt einen einfachen Befehl: \enquote{Heb den gestreiften Würfel an.}
Für einen Menschen ist das ziemlich simpel.
Ein KI-Agent muss dagegen zuerst wissen, wie die Szene aussieht, und diese Information muss ihm in einer Form vorliegen, die er auch verarbeiten kann.

Grundsätzlich gibt es dafür zwei Wege.
Die Szene lässt sich strukturell beschreiben, also als exakte Daten über die Objekte mit ihren Typen, Positionen und Größen, oder visuell als gerendertes Bild, das ein visuelles Modell interpretiert.
Beide Wege haben Stärken und blinde Flecken.
Die strukturellen Daten kennen jede Koordinate, über das Aussehen eines Objekts sagen sie nichts, und beim Bild ist es andersherum.

Der gestreifte Würfel aus dem Beispiel zeigt den ersten blinden Fleck.
Da das Streifenmuster in keinen strukturellen Daten steht, lässt sich der Würfel nur über das Bild bestimmen.
Ebenso gibt es den umgekehrten Fall.
Steht ein Objekt hinter einem anderen, taucht es im Bild überhaupt nicht auf, obwohl es in den strukturellen Daten vollständig verzeichnet ist.
Ein Befehl, der sich auf ein solches Objekt bezieht, lässt sich nur über die Struktur auflösen.

Wer einen solchen Agenten entwickelt, muss deshalb früh entscheiden, in welcher Form er ihm die Szene bereitstellt und ob sich beide Wege sinnvoll kombinieren lassen.
Besonders wichtig ist diese Entscheidung in immersiven Umgebungen, in denen Nutzer über natürliche Sprache mit einer Szene interagieren.
Dort bieten sich lokal betreibbare Modelle an, vor allem wegen Datenschutz, Latenz und Offline-Fähigkeit, und für solche Modelle gibt es bisher keinen systematischen Vergleich der beiden Wege.

Mit \emph{lokal} ist in dieser Arbeit durchgängig ein selbst gehosteter Betrieb auf eigener Hardware ohne Cloud-Dienst gemeint, bei dem keine Szenen- oder Sprachdaten das eigene Netz verlassen.
Auf dem Endgerät selbst müssen die Modelle dafür nicht rechnen.
Sowohl in der Messung als auch im Demonstrator laufen die Modelle auf einem Arbeitsplatzrechner, den die Brille im Demonstrator über eine Netzverbindung innerhalb des Institutsnetzes erreicht (Abschnitt~\ref{sec:impl-demonstrator}).


% ------------------------------------------------------------
\section{Relevanz}
\label{sec:ein-relevanz}

Die Wahl der Repräsentationsform entscheidet mit darüber, welche sprachlichen Anweisungen ein Agent überhaupt zuverlässig umsetzen kann.
Fehlt in einer Repräsentation eine ganze Klasse von Referenzen, scheitern die zugehörigen Befehle, egal wie gut der Agent ansonsten arbeitet.

Große, cloudbasierte Modelle können eine schwache Repräsentation teilweise durch ihre Größe ausgleichen.
Bei kleinen, lokal betriebenen Modellen ist das kaum möglich, weshalb die Repräsentation hier eine deutlich größere Rolle spielt.
Die Entscheidung über die Repräsentationsform fällt früh im Entwurf, hat weitreichende Folgen und wird bislang ohne belastbare Grundlage getroffen.

Darüber hinaus betrifft die Frage auch ein allgemeineres Problem.
Sobald ein Agent strukturierte und visuelle Information gemeinsam verarbeitet, stellt sich die Frage, ob sich beide Informationsarten ergänzen oder ob eine die andere verdrängt.
Auch über den hier untersuchten Anwendungsfall hinaus ist diese Frage noch offen~\cite{deng2025, wu2025}.


% ------------------------------------------------------------
\section{Zielsetzung und Forschungsfragen}
\label{sec:ein-ziele}

Diese Arbeit untersucht, welche Repräsentationsform einen lokal betriebenen Agenten bei der natürlichsprachlichen Szenenmanipulation am zuverlässigsten unterstützt.
Dafür werden zwei Forschungsfragen gestellt.

\begin{quote}
\textbf{Forschungsfrage~1 (RQ1).}
Welchen Einfluss hat die Form der Szenenrepräsentation, strukturell, visuell oder hybrid, auf den Ausführungserfolg der natürlichsprachlichen Szenenmanipulation durch einen lokal betriebenen Agenten, und zwar aufgeschlüsselt nach der Art der sprachlichen Referenz?

\vspace{6pt}
\textbf{Forschungsfrage~2 (RQ2).}
Lässt sich das Zusammenspiel der strukturellen und der visuellen Schicht allein durch die Gestaltung der Instruktion so steuern, dass die hybride Repräsentation die Stärken beider Schichten vereint?
\end{quote}

Die Stärken beider Schichten gelten für RQ2 dann als vereint, wenn eine gestaltete Instruktion dort deutlich besser abschneidet, wo die bloße Kombination beider Schichten scheitert, und in den übrigen Aufgaben nicht schlechter wird.

Gemessen wird in beiden Fällen ausschließlich der Ausführungserfolg.
Er wird für jeden Lauf binär aus dem tatsächlichen Szenenzustand bestimmt.
Ein Lauf ist nur erfolgreich, wenn das gemeinte Objekt betroffen ist und das geforderte Ergebnis eintritt, bei Positionsaufgaben innerhalb einer festen Toleranz.
Ein kontinuierlicher Lagefehler gehört nicht zur Zielgröße und wird nicht berichtet.
Die genaue Definition folgt in Abschnitt~\ref{sec:eval-metrics}.

Untersucht wird ein lokales Modellpaar in kontrolliert aufgebauten Szenen, gemessen wird am Desktop.
Zusätzlich entsteht ein Demonstrator, der die zuverlässigste Variante in die immersive Zielumgebung überträgt und die praktische Machbarkeit zeigen soll, ohne eine eigene Forschungsfrage zu beantworten.


% ------------------------------------------------------------
\section{Methodisches Vorgehen}
\label{sec:ein-vorgehen}

Die Arbeit verfolgt einen empirischen Ansatz.
Grundlage ist ein prototypisches System, das dieselbe Szene in allen drei Repräsentationsformen bereitstellt und über ein gemeinsames Interface manipulierbar macht.
Mit diesem System werden die Formen unter gleichen Bedingungen an kontrolliert aufgebauten Szenarien miteinander verglichen.
Ob eine Manipulation erfolgreich war, wird aus dem tatsächlichen Zustand der Szene bestimmt und nicht aus der Rückmeldung des Agenten, und jede Messung wird mehrfach wiederholt, um zufällige von belastbaren Unterschieden trennen zu können.

Daraus ergeben sich drei Beiträge, die in Abschnitt~\ref{sec:zus-beitrag} noch einmal aufgegriffen und eingeordnet werden.
Inhaltlich zeigt die Arbeit an einem kontrollierten Fall, dass sich strukturelle und visuelle Information nicht von selbst ergänzen und ihre Verbindung bewusst gestaltet werden muss.
Methodisch wird eine nach Referenztypen aufgeschlüsselte Auswertung eingeführt, die Zuständigkeiten der Repräsentationen sichtbar macht, die bei einem pauschalen Vergleich verborgen blieben.
Praktisch entsteht ein lauffähiger Demonstrator, der die gesamte Kette bis in die immersive Zielumgebung überträgt.


% ------------------------------------------------------------
\section{Aufbau der Arbeit}
\label{sec:ein-aufbau}

In Kapitel~\ref{ch:grundlagen} werden die für die Arbeit notwendigen Grundlagen dargestellt, die Arbeit in den Stand der Technik eingeordnet und die Forschungslücke abgeleitet.
Kapitel~\ref{ch:anforderungen} beschreibt, welche Anforderungen ein System erfüllen muss, um die Forschungsfragen beantworten zu können.
In Kapitel~\ref{ch:konzept} wird daraus das Lösungskonzept entwickelt, vor allem die drei Repräsentationsformen und die Entscheidungen, die einen fairen Vergleich ermöglichen.
Kapitel~\ref{ch:implementierung} beschreibt die technische Umsetzung, Kapitel~\ref{ch:evaluation} enthält den Vergleich und die Ergebnisse zu beiden Forschungsfragen.
Kapitel~\ref{ch:zusammenfassung} schließt mit einer Zusammenfassung und einem Ausblick.
